% GENERATED FILE. Edit canonical lessons, then run npm run generate:books.
% canonical-source-hash: fnv1a64:a2e4d5a01b586750
% canonical-lessons: JA-C131-ocha, JA-W131-small-ya, JA-C131-isha, JA-C131-chotto, JA-W131-small-yo, JA-C131-toshokan, JA-C131-ocha-o-kudasai, JA-W131-wo, JA-R131-tea-please

\chapter{Tea, Please: Small Ya, Small Yo and Wo}
\label{ch:ja-tea-please-small-ya-small-yo-wo}
\hlchaptermodality{131}

\begin{chapteropening}
\textbf{By the end of this chapter:} \emph{I can write small ya, small yo and wo, read words that use them, and ask for tea politely.}

Read ocha, isha, chotto and toshokan with their beats, write ocha o kudasai from memory with \ja{を} for the particle, and write the three new signs once each.
\end{chapteropening}

\section[\texorpdfstring{\ja{おちゃ} (ocha)}{おちゃ (ocha)}]{\ja{おちゃ} (ocha) — tea}

\label{lesson:JA-C131-ocha}



\begin{quote}
Before the new word, say \emph{o}, then \emph{chi}. \emph{Write it:} \textbf{\ja{お}}, then \textbf{\ja{ち}}

\begin{itemize}
\raggedright
  \item \emph{Recall:} say \emph{a sweetfish} — \textbf{R4}, eighty lessons back
  \item \emph{Recall:} say \emph{refreshing} — \textbf{R3}, twenty lessons back
\end{itemize}
\end{quote}

\subsection*{You'll want to know: \ja{おちゃ}}
\textbf{\ja{おちゃ}} — \emph{ocha} — \textquotedblleft{}tea\textquotedblright{}.

Count the beats as you say it: \emph{o–cha}. Two beats, not three. The last two signs are one beat between them, and the small sign at the end is the reason. The next lesson takes that small sign and writes it on its own.

\subsection*{Your turn}
\begin{itemize}
\raggedright
  \item \emph{Say it:} \emph{ocha}
  \item \emph{Say it:} \emph{ocha}, then count its beats aloud — two
  \item \emph{Recall:} name the sign in \textbf{\ja{おちゃ}} you can already write, and the one you cannot
\end{itemize}

\begin{tcolorbox}[breakable,colback=teal!4,colframe=teal!35!black,title={Before you move on}]
What is \textquotedblleft{}tea\textquotedblright{}? (\textbf{\ja{おちゃ}}, \emph{ocha}.) How many beats? (\textbf{Two}.)
\end{tcolorbox}
\section[\texorpdfstring{\ja{ゃ} (small ya)}{ゃ (small ya)}]{\ja{ゃ} — \ja{や}, written small}

\label{lesson:JA-W131-small-ya}



\begin{quote}
Five recalls, then the sign.

\begin{itemize}
\raggedright
  \item \emph{Recall:} say \emph{tea}, and how many beats it has — \textbf{R1}, one lesson back
  \item \emph{Recall:} write \textbf{\ja{や}}, full size: sweep, tick, descender
  \item \emph{Recall:} write \textbf{\ja{ゅ}}, and say what it does to the sign before it
  \item \emph{Recall:} say \emph{a trout} — \textbf{R4}, eighty lessons back
  \item \emph{Recall:} say \emph{calm} — \textbf{R3}, twenty lessons back
\end{itemize}
\end{quote}

\subsection*{Script — the shape}
\hlblockfigure{\detokenize{figures/JA-W131-small-ya-filmstrip.pdf}}{How \ja{ゃ} is written, stroke by stroke}

\begin{quote}
\textbf{\ja{ゃ}} — small \emph{ya}
\end{quote}

Same shape as \textbf{\ja{や}}, in the same order, written small and sitting low: the sweep and its hook first, then the short tick, then the descender. Three strokes, two lifts.

It does the job \textbf{\ja{ゅ}} does. It joins to the sign before it, and the two become \textbf{one beat}:

\noindent
\begin{tabularx}{\linewidth}{@{}>{\raggedright\arraybackslash}X>{\raggedright\arraybackslash}X>{\raggedright\arraybackslash}X@{}}
\toprule
\textbf{you write} & \textbf{you say} & \textbf{beats} \\
\midrule
\textbf{\ja{ちや}} & \emph{chi-ya} & two \\
\textbf{\ja{ちゃ}} & \emph{cha} & one \\
\bottomrule
\end{tabularx}

That is the whole difference between the two rows: the size of one sign. The tea word you met one lesson back is \emph{o–cha}, two beats, because its \textbf{\ja{ちゃ}} is one.

\subsection*{Your turn}
\begin{enumerate}
\raggedright
  \item Write \textbf{\ja{や}}, then \textbf{\ja{ゃ}} beneath it, half the height, in the same three strokes.
  \item Write \textbf{\ja{おちゃ}} from memory. Check that the last sign sits low and small.
  \item Read \textbf{\ja{ちや}} and \textbf{\ja{ちゃ}} aloud, one after the other, and clap the beats.
\end{enumerate}

\begin{tcolorbox}[breakable,colback=teal!4,colframe=teal!35!black,title={Before you move on}]
Stop after one accurate small form.

Source: \href{https://www.unicode.org/charts/PDF/U3040.pdf}{Unicode Hiragana chart}; stroke count and order per \href{https://kanjivg.tagaini.net/kanjivg/kanji/03083.html}{KanjiVG}.
\end{tcolorbox}
\section[\texorpdfstring{\ja{いしゃ} (isha)}{いしゃ (isha)}]{\ja{いしゃ} (isha) — a doctor}

\label{lesson:JA-C131-isha}



\begin{quote}
\emph{Write it:} \textbf{\ja{ゃ}} — \textbf{R1}, one lesson back

Then say \emph{ocha}.

\begin{itemize}
\raggedright
  \item \emph{Recall:} say \emph{a crucian carp} — \textbf{R4}, eighty lessons back
  \item \emph{Recall:} say \emph{honest} — \textbf{R3}, twenty lessons back
  \item \emph{Recall:} say \emph{bright red} — \textbf{R2}, seven lessons back
\end{itemize}
\end{quote}

\subsection*{You'll want to know: \ja{いしゃ}}
\textbf{\ja{いしゃ}} — \emph{isha} — \textquotedblleft{}a doctor\textquotedblright{}.

The small sign is spent at once. \textbf{\ja{し}} with \textbf{\ja{ゃ}} after it is one beat, \emph{sha}, so the word is two beats: \emph{i–sha}.

\subsection*{Your turn}
\begin{itemize}
\raggedright
  \item \emph{Say it:} \emph{isha}
  \item \emph{Say it:} \emph{isha}, then count its beats aloud — two
  \item \emph{Write it:} \textbf{\ja{いしゃ}} from memory, keeping the last sign small
\end{itemize}

\begin{tcolorbox}[breakable,colback=teal!4,colframe=teal!35!black,title={Before you move on}]
What does \ja{いしゃ} mean? (\textquotedblleft{}A doctor\textquotedblright{}.) Say it once more.
\end{tcolorbox}
\section[\texorpdfstring{\ja{ちょっと} (chotto)}{ちょっと (chotto)}]{\ja{ちょっと} (chotto) — a little}

\label{lesson:JA-C131-chotto}



\begin{quote}
Say \emph{a doctor} — \textbf{R1}, one lesson back. \emph{Write it:} the small sign in it

\begin{itemize}
\raggedright
  \item \emph{Recall:} say \emph{a young yellowtail} — \textbf{R4}, eighty lessons back
  \item \emph{Recall:} say \emph{clever} — \textbf{R3}, twenty lessons back
  \item \emph{Recall:} say \emph{vague} — \textbf{R2}, seven lessons back
\end{itemize}
\end{quote}

\subsection*{You'll want to know: \ja{ちょっと}}
\textbf{\ja{ちょっと}} — \emph{chotto} — \textquotedblleft{}a little\textquotedblright{}, and, said on its own with a small pause, \textquotedblleft{}one moment\textquotedblright{}.

Three beats: \emph{cho}, a held silence, \emph{to}. The silence is the small \textbf{\ja{っ}} you already write. The first beat is \textbf{\ja{ち}} with a small sign after it that you do not write yet; it joins forward, the way \textbf{\ja{ゃ}} does, and the next lesson writes it.

\subsection*{Your turn}
\begin{itemize}
\raggedright
  \item \emph{Say it:} \emph{chotto}
  \item \emph{Say it:} \emph{chotto}, then count its beats aloud — three, and the middle one is a silent hold
  \item \emph{Recall:} name the two small signs in \textbf{\ja{ちょっと}} and say which one holds and which one joins
\end{itemize}

\begin{tcolorbox}[breakable,colback=teal!4,colframe=teal!35!black,title={Before you move on}]
What is \textquotedblleft{}a little\textquotedblright{}? (\textbf{\ja{ちょっと}}, \emph{chotto}.) How many beats? (\textbf{Three}.)
\end{tcolorbox}
\section[\texorpdfstring{\ja{ょ} (small yo)}{ょ (small yo)}]{\ja{ょ} — \ja{よ}, written small}

\label{lesson:JA-W131-small-yo}



\begin{quote}
Six recalls, then the sign.

\begin{itemize}
\raggedright
  \item \emph{Recall:} say \emph{a little}, and clap its three beats — \textbf{R1}, one lesson back
  \item \emph{Recall:} say \emph{a doctor} — \textbf{R1}, two lessons back
  \item \emph{Recall:} write \textbf{\ja{よ}}, full size: the bar, then the stem and its loop
  \item \emph{Recall:} say \emph{a Pacific saury} — \textbf{R4}, eighty lessons back
  \item \emph{Recall:} say \emph{foolish} — \textbf{R3}, twenty lessons back
  \item \emph{Recall:} say \emph{certain} — \textbf{R2}, seven lessons back
\end{itemize}
\end{quote}

\subsection*{Script — the shape}
\hlblockfigure{\detokenize{figures/JA-W131-small-yo-filmstrip.pdf}}{How \ja{ょ} is written, stroke by stroke}

\begin{quote}
\textbf{\ja{ょ}} — small \emph{yo}
\end{quote}

Same shape as \textbf{\ja{よ}}, in the same order, written small and sitting low: the short bar from left to right, then the stem down through it and round the loop. Two strokes, one lift.

You now have all three small signs that join forward. Each one turns the sign before it into a single beat:

\noindent
\begin{tabularx}{\linewidth}{@{}>{\raggedright\arraybackslash}X>{\raggedright\arraybackslash}X@{}}
\toprule
\textbf{you write} & \textbf{you say, one beat each} \\
\midrule
\textbf{\ja{ちゃ}}, \textbf{\ja{ちゅ}}, \textbf{\ja{ちょ}} & \emph{cha}, \emph{chu}, \emph{cho} \\
\textbf{\ja{しゃ}}, \textbf{\ja{しゅ}}, \textbf{\ja{しょ}} & \emph{sha}, \emph{shu}, \emph{sho} \\
\bottomrule
\end{tabularx}

So \textbf{\ja{ちょ}} is one beat, \emph{cho}, and \textbf{\ja{ちょっと}} is the three beats you said in the last lesson: \emph{cho}, a held silence, \emph{to}.

\subsection*{Your turn}
\begin{enumerate}
\raggedright
  \item Write \textbf{\ja{よ}}, then \textbf{\ja{ょ}} beneath it, half the height.
  \item Write \textbf{\ja{ちょっと}} from memory: one small sign that joins, one that holds.
  \item Write \textbf{\ja{しゃ}}, \textbf{\ja{しゅ}}, \textbf{\ja{しょ}} in a row and read them as three single beats.
\end{enumerate}

\begin{tcolorbox}[breakable,colback=teal!4,colframe=teal!35!black,title={Before you move on}]
Stop after one accurate small form.

Source: \href{https://www.unicode.org/charts/PDF/U3040.pdf}{Unicode Hiragana chart}; stroke count and order per \href{https://kanjivg.tagaini.net/kanjivg/kanji/03087.html}{KanjiVG}.
\end{tcolorbox}
\section[\texorpdfstring{\ja{としょかん} (toshokan)}{としょかん (toshokan)}]{\ja{としょかん} (toshokan) — a library}

\label{lesson:JA-C131-toshokan}



\begin{quote}
\emph{Write it:} \textbf{\ja{ょ}} — \textbf{R1}, one lesson back

Then say \emph{tea} — \textbf{R2}, five lessons back.

\begin{itemize}
\raggedright
  \item \emph{Recall:} say \emph{a short-neck clam} — \textbf{R4}, eighty lessons back
  \item \emph{Recall:} say \emph{brave} — \textbf{R3}, twenty lessons back
\end{itemize}
\end{quote}

\subsection*{You'll want to know: \ja{としょかん}}
\textbf{\ja{としょかん}} — \emph{toshokan} — \textquotedblleft{}a library\textquotedblright{}.

Five signs, four beats: \emph{to–sho–ka–n}. \textbf{\ja{し}} with \textbf{\ja{ょ}} after it is one beat, \emph{sho}, and \textbf{\ja{ん}} is a beat on its own.

\subsection*{Your turn}
\begin{itemize}
\raggedright
  \item \emph{Say it:} \emph{toshokan}
  \item \emph{Say it:} \emph{toshokan}, then count its beats aloud — four
  \item \emph{Write it:} \textbf{\ja{としょかん}} from memory, keeping the second sign small
\end{itemize}

\begin{tcolorbox}[breakable,colback=teal!4,colframe=teal!35!black,title={Before you move on}]
What does \ja{としょかん} mean? (\textquotedblleft{}A library\textquotedblright{}.) How many beats? (\textbf{Four}.)
\end{tcolorbox}
\section[\texorpdfstring{\ja{おちゃを} \ja{ください} (ocha o kudasai)}{おちゃを ください (ocha o kudasai)}]{\ja{おちゃを} \ja{ください} (ocha o kudasai) — tea, please}

\label{lesson:JA-C131-ocha-o-kudasai}



\begin{quote}
Say \emph{a library} — \textbf{R1}, one lesson back. Then say \emph{please say it}. \emph{Write it:} \textbf{\ja{ゃ}} — \textbf{R2}, five lessons back

\begin{itemize}
\raggedright
  \item \emph{Recall:} say \emph{a scallop} — \textbf{R4}, eighty lessons back
  \item \emph{Recall:} say \emph{fine, small and detailed} — \textbf{R3}, twenty lessons back
\end{itemize}
\end{quote}

\subsection*{You'll want to know: \ja{おちゃを} \ja{ください}}
\textbf{\ja{おちゃを} \ja{ください}} — \emph{ocha o kudasai} — \textquotedblleft{}tea, please\textquotedblright{}.

You know both words already: \emph{ocha}, tea, and \emph{kudasai}, the polite request at the end of \emph{itte kudasai}. What is new is the small word between them, said \emph{o}.

It marks \textbf{the thing the request is about}: tea is what you are asking for. Put any thing you can name in front of it and the request works the same way.

That word is spelled with a sign you have not written yet. It is said \emph{o}, the same as \textbf{\ja{お}}, but it is a different sign, and the next lesson writes it.

\subsection*{Your turn}
\begin{itemize}
\raggedright
  \item \emph{Say it:} \emph{ocha o kudasai}
  \item \emph{Say it:} \emph{ocha o kudasai}, with a short pause after \emph{o}
  \item \emph{Recall:} say which of the three words names the thing, which one marks it, and which one asks
\end{itemize}

\begin{tcolorbox}[breakable,colback=teal!4,colframe=teal!35!black,title={Before you move on}]
How do you ask for tea? (\textbf{\ja{おちゃを} \ja{ください}}, \emph{ocha o kudasai}.)
\end{tcolorbox}
\section[\texorpdfstring{\ja{を} (wo)}{を (wo)}]{\ja{を} — said o, kept for one job}

\label{lesson:JA-W131-wo}



\begin{quote}
Five recalls, then the sign.

\begin{itemize}
\raggedright
  \item \emph{Recall:} ask for tea politely — \textbf{R1}, one lesson back
  \item \emph{Recall:} say \emph{a doctor} — \textbf{R2}, five lessons back
  \item \emph{Recall:} write \textbf{\ja{お}}, and say it
  \item \emph{Recall:} say \emph{a flea} — \textbf{R4}, eighty lessons back
  \item \emph{Recall:} say \emph{ugly} — \textbf{R3}, twenty lessons back
\end{itemize}
\end{quote}

\subsection*{Script — the shape}
\hlblockfigure{\detokenize{figures/JA-W131-wo-filmstrip.pdf}}{How \ja{を} is written, stroke by stroke}

\begin{quote}
\textbf{\ja{を}} — said \emph{o}
\end{quote}

Three strokes, two lifts:

\begin{enumerate}
\raggedright
  \item the short \textbf{bar} at the top, left to right
  \item starting above the bar, a long \textbf{diagonal} down and to the left through it; at its foot, without lifting, turn sharply back up and to the right into an \textbf{arch}, and come down its right side in a short stem
  \item a separate \textbf{curve}: start at the right, sweep down and left across the middle, round the lower left, and finish along the base to the right
\end{enumerate}

It is said exactly like \textbf{\ja{お}}. The two are not interchangeable, and what separates them is not sound:

\begin{itemize}
\raggedright
  \item \textbf{\ja{お}} spells \emph{o} inside words, as in \textbf{\ja{おちゃ}}.
  \item \textbf{\ja{を}} spells the particle that marks the thing asked for or acted on, and nothing else.
\end{itemize}

You have met this once already. \textbf{\ja{は}} is read \emph{wa} when it is the topic marker; \textbf{\ja{を}} is read \emph{o} when it is the object marker. Both are spellings kept for a grammatical word after the sound moved.

\subsection*{Your turn}
\begin{enumerate}
\raggedright
  \item Trace \textbf{\ja{を}} saying \textquotedblleft{}bar, diagonal and arch, curve\textquotedblright{}.
  \item Write \textbf{\ja{お}} and \textbf{\ja{を}} side by side; say both, and hear that they are the same.
  \item Write \textbf{\ja{おちゃを} \ja{ください}} from memory. Check which \emph{o} is which.
\end{enumerate}

\begin{tcolorbox}[breakable,colback=teal!4,colframe=teal!35!black,title={Before you move on}]
Which sign spells the particle? (\textbf{\ja{を}}.)

Source: \href{https://www.unicode.org/charts/PDF/U3040.pdf}{Unicode Hiragana chart}; stroke count and order per \href{https://kanjivg.tagaini.net/kanjivg/kanji/03092.html}{KanjiVG}.
\end{tcolorbox}
\section[(dialogue)]{Tea, please — the three signs, written}

\label{lesson:JA-R131-tea-please}



\begin{quote}
\emph{Write it:} \textbf{\ja{を}} — \textbf{R1}, one lesson back

Then say \emph{a little} — \textbf{R2}, five lessons back.

\begin{itemize}
\raggedright
  \item \emph{Recall:} say \emph{a bat} — \textbf{R4}, eighty lessons back
  \item \emph{Recall:} say \emph{cute} — \textbf{R3}, twenty lessons back
\end{itemize}
\end{quote}

\subsection*{Your turn — read the four words}
\emph{Read it:} each one below aloud, count its beats, then say what it means

Read these aloud:

\begin{itemize}
\raggedright
  \item \textbf{\ja{おちゃ}} — two beats
  \item \textbf{\ja{いしゃ}} — two beats
  \item \textbf{\ja{ちょっと}} — three beats
  \item \textbf{\ja{としょかん}} — four beats
\end{itemize}

Every small sign that joins forward folds two signs into one beat. The small \textbf{\ja{っ}} is the only small sign here that takes a beat of its own.

\subsection*{Your turn — write the request}
\begin{enumerate}
\raggedright
  \item \emph{Write it:} \textbf{\ja{おちゃを} \ja{ください}} from memory, and circle the two signs said \emph{o}
  \item \emph{Write it:} \textbf{\ja{ゃ}}, \textbf{\ja{ょ}} and \textbf{\ja{を}} once each — small signs low and half height
  \item Say \emph{ocha o kudasai} to an imaginary waiter, then \emph{chotto}, as if asking them to wait a moment.
\end{enumerate}

\begin{tcolorbox}[breakable,colback=teal!4,colframe=teal!35!black,title={Before you move on}]
Tea? (\textbf{\ja{おちゃ}}, \emph{ocha}.) A doctor? (\textbf{\ja{いしゃ}}, \emph{isha}.) A little? (\textbf{\ja{ちょっと}}, \emph{chotto}.) A library? (\textbf{\ja{としょかん}}, \emph{toshokan}.) Tea, please? (\textbf{\ja{おちゃを} \ja{ください}}, \emph{ocha o kudasai}.)
\end{tcolorbox}
